% PDF 1: 封面。
\definecolor{coverblue}{HTML}{214358}
\definecolor{coverlight}{HTML}{A4D2E8}
\pagecolor{coverblue}
\begin{titlepage}
\color{white}
\pdfbookmark[0]{封面}{cover}
\noindent\includegraphics[width=\linewidth]{assets/cover_art.pdf}\par
\vspace{21pt}
\noindent{\sffamily\large\color{coverlight}Uday Kamath · Kevin Keenan\par
\noindent Garrett Somers · Sarah Sorenson\par}
\vspace{22pt}
\noindent{\sffamily\bfseries\fontsize{28}{36}\selectfont 大语言模型深度解析\par}
\vspace{10pt}
\noindent{\sffamily\Large\color{coverlight}贯通理论与实践\par}
\vspace{16pt}
\noindent{\large Large Language Models: A Deep Dive\par}
\noindent Bridging Theory and Practice\par
\vfill
\iftranslationcomplete\else
\noindent{\sffamily 译排进行中}\par\medskip
\fi
\noindent 原版由 Springer 出版，2024\par
\end{titlepage}
\nopagecolor

% PDF 2--3: 半书名页与书名页。
\thispagestyle{empty}
\vspace*{15mm}
\noindent{\sffamily\bfseries\Large 大语言模型深度解析\par}
\clearpage
\thispagestyle{empty}
\vspace*{15mm}
\noindent{\large Uday Kamath · Kevin Keenan\par
\noindent Garrett Somers · Sarah Sorenson\par}
\vspace{23mm}
\noindent{\sffamily\bfseries\Huge 大语言模型深度解析\par}
\medskip
\noindent{\sffamily\Large 贯通理论与实践\par}
\vfill
\noindent Springer · Cham
\clearpage

% PDF 4: 出版和版权信息。
\thispagestyle{empty}
\noindent
\begin{tabularx}{\linewidth}{@{}XX@{}}
Uday Kamath & Kevin Keenan\\
Smarsh & Smarsh\\
美国弗吉尼亚州阿什本 & 英国唐郡卡斯尔韦兰\\[8pt]
Garrett Somers & Sarah Sorenson\\
Smarsh & Smarsh\\
美国田纳西州纳什维尔 & 美国弗吉尼亚州雷斯顿
\end{tabularx}
\vfill
\begingroup
\small
\setlength{\parindent}{0pt}
ISBN 978-3-031-65646-0\par
ISBN 978-3-031-65647-7（电子书）\par
\url{https://doi.org/10.1007/978-3-031-65647-7}\par\smallskip

\textcopyright\ 编者（如有）及作者，2024 年独家许可予 Springer Nature Switzerland AG。

本书受著作权保护。无论涉及全部材料还是其中一部分，所有权利均由出版方独占取得许可。这些权利尤其包括翻译、重印、再使用插图、朗诵、广播、以微缩胶片或其他实物方式复制、传输或存储和检索信息、电子改编、以计算机软件使用，以及采用任何已知或日后发展出的类似或不同方法使用的权利。

本书使用一般性描述名称、注册名称、商标、服务标记等。即使未作专门说明，也不能据此认定这些名称不受有关保护性法律法规约束，可以自由使用。

出版方、作者和编者认为，本书所载建议和信息在出版时真实、准确。出版方、作者和编者均不就本书所载材料或其中可能存在的错误或遗漏作出任何明示或默示保证。对于所刊地图和机构隶属关系中涉及的管辖权主张，出版方保持中立。\par\smallskip

本书以 Springer 品牌出版，出版公司为 Springer Nature Switzerland AG。\par
公司注册地址为 Gewerbestrasse 11, 6330 Cham, Switzerland。\par\smallskip
若需处置本书，请将纸张回收。
\endgroup
\clearpage

% PDF 5: 献词。
\thispagestyle{empty}
\vspace*{28mm}
\begin{flushright}
献给我的父母 Krishna 和 Bharathi、妻子 Pratibha、孩子 Aaroh 和 Brandy，以及我亲爱的家人和朋友 Anu、Jyoti、Maithilee、Munnu、Nana、Priyal、Shilpa、Sunila、Swati 和 Swetha。感谢你们始终如一的支持、爱和关怀，尤其是在我患病期间。\par
——Uday Kamath\par\bigskip

献给我的父母 Eugene 和 Evelyn，始终耐心、深爱着我的妻子 Kaitlin，以及我美丽的孩子 Éirinn 和 Méabh。你们欣然支持我做过许多事，而其中少有像写这本书这样雄心勃勃的事。\par
——Kevin Keenan\par\bigskip

献给我出色的女儿们 Octavia 和 Lyra。\par
——Garrett Somers\par\bigskip

献给我的丈夫 Matthew。写作期间，我常常抽不出时间，他一一体谅，只有偶尔才抱怨几句。\par
——Sarah Sorenson
\end{flushright}

% PDF 6--7: 序。
\frontchapter{序}
大语言模型（Large Language Model，LLM）彻底改变了人工智能（Artificial Intelligence，AI）领域，正在改变人们与技术交互的方式，也在重塑多个行业。我在牛津大学（University of Oxford）担任多门人工智能课程的课程主任，同时参与世界各地多个创业项目，亲眼见到这些模型如何解决复杂问题、简化日常工作。本书面世正逢其时，为理解和使用大语言模型提供了全面的指引。作者把模型背后的复杂架构和算法逐步拆解，讲得明白，使不同背景的读者都能理解。

我认识第一作者 Uday Kamath 博士已有一段时间，也一直关注他的工作。他在人工智能领域的专长和见解广受认可。本书体现了他对这一领域的深入理解和创新思考。从预训练（pre-training）与基于提示的学习（prompt-based learning）的基础，到微调（fine-tuning）方法和检索增强生成（Retrieval-Augmented Generation，RAG）等进阶主题，书中都作了详细介绍。我尤其看重它对实践的重视。每章都着眼于为读者提供在真实场景中应用大语言模型所需的知识和技能。通过动手教程和实际案例，读者可以理解理论，也会更有把握在工作中有效地应用这些模型。

专门讲解大语言模型运维（Large Language Model Operations，LLMOps）和生产部署（productionizing）的一章尤其有价值。它细致说明怎样在实际环境中部署和运行这些模型，让读者把理论知识转化为实际成果。书中还汇集了大量数据集、基准和评估指标，可供读者探索大语言模型应用。多模态大语言模型（Multimodal Large Language Model，MMLLM）一章则将视野从文本扩展到音频、图像、视频以及机器人的输入和输出，展现了人工智能交互的新方向。本书也讨论了与大语言模型相关的伦理难题，包括偏见（bias）、公平性（fairness）和隐私。在发挥这些模型能力的同时，负责任地使用它们至关重要。作者提出了应对这些问题的策略，有助于使模型的使用符合人类价值观。

总之，凡是对大语言模型感兴趣的读者，都可从本书获益。它深入讲解技术，提供实践经验，并讨论伦理问题，是了解人工智能未来发展的指南。作者考察详实，表述清楚。我相信，本书在未来多年都将是这一领域的重要参考。
\reviewer{——Ajit Jaokar，牛津大学人工智能课程主任、访问研究员}

% PDF 8--9: 推荐语。
\frontchapter{推荐语}
“作为一名资深技术高管和行业意见领袖，我意识到，大语言模型在产品开发各个环节中的作用越来越重要，从提高开发效率到推动产品和营销落地，无一例外。为了打牢基础，我读了这本大语言模型著作，发现它是一部难得的指南。在我和团队探索如何把大语言模型的能力用于产品时，它已经成了我随时查阅的参考书。

“本书揭开了大语言模型的复杂面纱，从精细的架构讲到部署中的伦理考量，也着眼于多模态大语言模型的未来，说明模型怎样从文本扩展到音频、图像和视频。在今天这个由人工智能驱动的世界，任何想理解并运用大语言模型变革力量的人，都值得读一读这本内容全面、富有实践见解的书。”
\reviewer{——Shalini Govil Pai，Google 副总裁兼总经理}

“文字准确，技术性强，适合具有扎实机器学习（machine learning）和人工智能基础的读者。全书信息密集，作者假定读者熟悉较高阶的数学概念并具备编程能力。尽管技术内容深入，本书结构分明，解释清楚，论述循序展开。”
\reviewer{——Sherry Marcus 博士，Amazon AWS 生成式人工智能应用科学总监}

“生成式人工智能（Generative AI）是当下的热门话题。它究竟会改变社会与商业的面貌，还是仅仅又一个流行词？要理解生成式人工智能的巨大潜力、风险和挑战，这本书值得一读。它完整走过生成式人工智能的生命周期，非常适合寻找全面指南的读者。开头用语言模型的发展史交代背景，接着深入讲解 Transformer 架构、提示工程（prompt engineering）、微调与检索增强生成，最后介绍部署这些方案的实用步骤。我觉得本书用途广泛，读来引人入胜，开发人员和人工智能爱好者都适合读。末章对生成式人工智能未来的展望尤其有见地。我诚挚推荐。”
\reviewer{——Eduardo Ordax，Amazon 生成式人工智能负责人、AWS 技术布道师}

“本书的覆盖面和时效性令人印象深刻。作者细致梳理了大语言模型近年的发展，将各种技术放到实际情境中组织、讲解。因此，对专业人士和研究人员而言，这是一份重要的参考，尤其适用于改善信息系统中的用户体验和交互。”
\reviewer{——Prithvi Prabhu，H2O.ai 首席技术专家}

“Kamath 与合著者为想深入理解大语言模型的研究人员和学生提供了一份全面的资料。书中巧妙连接理论和实践，也纳入多模态模型等前沿进展，因此对学术研究人员和产业界科学家都很有吸引力，值得常备。”
\reviewer{——Amarda Shehu 博士，乔治梅森大学（George Mason University）计算机科学教授，工程与计算学院分管人工智能创新的副院长}

“本书概述了构建可靠的大语言模型软件所需的知识。即使数据相当复杂，你也可以了解如何围绕产品和业务目标选择、调整、评估模型，以及如何借助 RAG 补充大语言模型系统。推荐所有人工智能工程师阅读！”
\reviewer{——Daniel Svonava，Superlinked 创始人兼首席执行官}

% PDF 10--17: 前言及原书两个清单。
\frontchapter{前言}
\frontsection{为什么写这本书}
在技术演进的长河中，大语言模型已经成为一块基石。它们正在改变我们与信息打交道的方式，重塑各个行业，并重新划定人工智能的边界。在这场转变初现之际，大语言模型的影响已经越过计算能力的进步本身，遍及从日常事务到复杂问题求解的各类活动。这一巨大变化既引起技术人员和研究者的兴趣，也让更多人想了解并运用大语言模型的能力。

大语言模型的出现标志着计算方式显著偏离了传统范式。机器能够以前所未有的方式理解、生成并借助人类语言互动，交互显得自然且富有见解。这类模型从大量人类语言数据中学习，由此为创新、创作和提高效率开辟了新的途径。它们处理和生成语言的能力带来了改变各领域的应用，包括内容自动生成、更复杂的聊天机器人、功能更强的搜索引擎，以及开创性的研究工具。

强大的能力伴随着复杂性。大语言模型的工作原理令人着迷，却并非人人都能立即理解。复杂的架构、底层算法，以及伴随模型部署的伦理问题，都十分重要，值得深入探讨。人们迫切需要一份全面的资料，既能解释这些模型在技术上如何运作，也能把影响放在具体背景下加以理解，探究应用，并讨论由此产生的伦理难题。这正是本书的目标。

本书深入考察大语言模型的实际工作方式。开篇综述预训练模型，按不同标准分类，并详解 Transformer 等架构。这为进一步理解基于提示的学习打下基础。书中介绍多种提示学习机制，说明从大语言模型提取知识的重要性及可用的技术，随后详细讨论微调的不同策略和取舍。

接下来，本书系统考察如何把强化学习（reinforcement learning，RL）纳入大语言模型的训练，让模型符合人类价值观。这是大语言模型训练的重要组成部分。我们还讨论大语言模型与计算机视觉（computer vision）、机器人学（robotics）、语音处理（speech processing）等领域的交汇。本书重视实际应用，详细介绍检索增强生成等真实用例。选择这些例子，是为了说明大语言模型在不同的行业和情境中可以发挥哪些作用。

本书也介绍如何使大语言模型投入运行和部署，引导读者借助现有工具和程序库实现这些模型，掌握实际使用方法。书中探讨模型所面临的挑战，从模型内在的偏见与不可预测性，到涌现行为（emergent behavior）所带来的更广泛伦理问题。最后，本书介绍前沿的多模态大语言模型。它们将处理范围从文本扩展到音频、图像、视频和机器人的输入，开启更全面的人工智能交互。各章还配有动手教程，展示大语言模型在自然语言处理（Natural Language Processing，NLP）任务中与本章主题相关的功能。

为了让读者能够深入而顺畅地认识大语言模型，本书提供以下内容。

\begin{itemize}
\item 探讨 100 多种技术和先进方法，包括预训练、基于提示的调优（prompt-based tuning）、指令调优（instruction tuning）、参数高效微调（Parameter-Efficient Fine-Tuning，PEFT）、计算高效微调（compute-efficient fine-tuning）、面向终端用户的提示工程、检索增强生成系统的构建与优化，以及使用强化学习使大语言模型符合人类价值观的策略。
\item 汇集 200 多个数据集，适用范围从预训练到多模态调优，可作为开发多种大语言模型应用的基础。
\item 介绍 50 多种应对幻觉（hallucination）、毒性（toxicity）、偏见、公平性和隐私等伦理问题的策略，说明如何度量、评估和缓解这些问题，以负责任地部署大语言模型。
\item 讲解如何选择合适的大语言模型、优化训练和推理（inference）成本、降低延迟、改善性能并利用相关工具，从而将模型的潜力付诸实践。
\item 收录 200 多项基准，覆盖各种任务上的模型表现、伦理考量和多模态应用，并介绍大语言模型生命周期中的 50 多项评估指标。
\item 通过 9 个详细教程，使用与 Google Colab 兼容的工具和程序库，练习大语言模型的预训练、微调、对齐调优（alignment tuning）、偏见缓解、多模态训练和部署，把理论概念用于实践。
\item 面向数据科学家和从业者提供 100 多条实用提示，介绍实现细节、技巧和工具，帮助读者处理大语言模型生命周期各环节的工作，高效完成任务。
\end{itemize}

\frontsection{本书适合谁读}
本书面向不同背景的读者，旨在为希望理解大语言模型基本思想和技术细节的人提供一站式参考。无论你是希望了解人工智能复杂机制的计算机科学、数据科学或人工智能专业本科生或研究生，是深入研究语言模型理论进展和实际应用的人工智能或自然语言处理研究者，还是想借助大语言模型的前沿能力解决实际问题的产业界数据科学家，本书都为你而写。

本书的内容安排兼顾初涉人工智能的读者，以及希望深化理解、扩展知识的资深从业者和学者。学生可以把它当作基础教材，学习理解和使用大语言模型所需的基本概念、方法和工具。它把课堂所学与实际工作中应对人工智能带来的挑战、把握机会所需的技能连接起来。

对于研究人员和学者，本书全面介绍大语言模型各方面的前沿研究，包括提示工程技术、学习策略、基于人类反馈的强化学习（Reinforcement Learning from Human Feedback，RLHF）和多模态大语言模型，也深入分析现有挑战及其缓解策略。

数据科学家和行业从业者可以用本书掌握高效微调大语言模型的技术，使模型适用于特定领域。除微调外，本书还讨论检索增强生成等应用与学习策略，帮助读者掌握把大语言模型部署到生产系统的技能。本书也深入讲解如何评估和缓解幻觉、偏见、公平性和隐私问题，使读者能够在实际环境中处理这些问题。

阅读本书之前，读者应具备一定的预备知识。
\begin{itemize}
\item 了解线性代数、微积分、统计学和概率论的基础。
\item 具备中等水平的机器学习和人工智能知识。
\item 了解基本的自然语言处理概念和深度学习技术。
\item 具备中等水平的 Python 编程能力，并熟悉 Pandas、Matplotlib、PyTorch 等程序库。
\end{itemize}

\frontsection{本书讲什么}
下面逐章概述内容与主题，让读者对将要学习的材料有更具体的认识。

\noindent\textbf{第 1 章\quad 大语言模型导论}\enspace 本章从自然语言处理的历史背景和演进谈起。它追溯人类语言能力的起源，说明如何逐步过渡到计算语言建模，突出其中生物演化与技术之间复杂联系的重要性。本章沿时间脉络介绍计算领域中语言模型的演进，展示早期的简单模型怎样发展为今天的大语言模型，讨论算法进展、算力和数据可获得性等促成变化的因素，并界定何谓大语言模型，帮助读者理解它在当代人工智能范式中的地位。

\noindent\textbf{第 2 章\quad 语言模型预训练}\enspace 本章深入介绍预训练模型及其基本机制和结构。首先细致考察注意力机制（attention mechanism），说明模型如何聚焦相关信息，以及这如何改变了自然语言处理。接着介绍 Transformer 架构的开创性及其在现代自然语言处理中的意义。本章随后对大语言模型分类，说明编码器—解码器（encoder-decoder）、自回归（autoregressive）语言模型和掩码语言模型（masked language model）的具体特点。讨论 BERT、T5、GPT（第 1—3 代）和 Mixtral~8x7B 等开创性架构时，侧重各自的训练技术和主要应用。随后介绍关键数据集，说明先进模型以哪些数据为基础。章末给出主要模型、工具和模型平台的实用指南，为后续进阶内容做好准备。

\noindent\textbf{第 3 章\quad 基于提示的学习}\enspace 本章探讨目前对自然语言处理发展十分重要的提示学习技术。它循序介绍基本原理，说明如何把各类自然语言处理任务映射为具体提示。然后讨论提示工程、答案工程（answer engineering）和多提示（multi-prompting），解释如何设计有效且高效的提示，引导模型给出所需输出。本章还将传统的预训练再微调方法与基于提示的方法作比较。

\noindent\textbf{第 4 章\quad 大语言模型的适配与使用}\enspace 本章研究大语言模型的能力及其在实践中带来的复杂问题，特别是必须微调模型的情形。书中详细考察多种参数高效学习策略，介绍串行适配器（serial adapter）、并行适配器（parallel adapter）、低秩适配（Low-Rank Adaptation，LoRA）和基于向量的随机矩阵适配（Vector-Based Random Matrix Adaptation，VeRA）等方法，并用增量调优（delta-tuning）的视角比较为达到预期效果而采用这些方法的效率。对数据有限的场景，本章阐述大语言模型的零样本、少样本和多样本学习（zero-, few-, and multi-shot learning）方法，说明提示设计和上下文长度的细节如何影响上下文学习（In-Context Learning，ICL），并强调数据稀少时思维链推理（chain-of-thought reasoning）的价值。最后，本章讨论如何降低全参数调优（full-parameter tuning）的资金和计算负担，介绍训练后量化（Post-Training Quantization，PTQ）和量化感知微调（quantization-aware fine-tuning）等做法，使资金有限的机构也有机会运用大语言模型的能力。

\noindent\textbf{第 5 章\quad 大语言模型的对齐调优}\enspace 本章介绍与人类偏好对齐的概念。人类偏好在此概括为 3H，即有帮助、无害和诚实（Helpful, Harmless, and Honest）。本章讨论把复杂的人类价值观写入大语言模型所遇到的难题，考察如何通过反馈机制，运用强化学习，尤其是基于人类反馈的强化学习，使大语言模型符合人类价值观。面对 RLHF 的资源需求高、难以扩展等挑战，本章介绍宪法式人工智能（Constitutional AI）和直接偏好优化（Direct Preference Optimization，DPO）等进展，说明怎样更合乎伦理、更负责任地应用大语言模型。

\noindent\textbf{第 6 章\quad 大语言模型的挑战及应对方法}\enspace 本章探讨大语言模型固有的难题和伦理困境。概述其局限和挑战之后，首先分析模型以海量且多样的数据训练而引发的认识论问题。接着考察预训练模型内部隐含的道德规范，探问模型自身有哪些偏见，以及可能在无意间传播哪些社会文化价值观。随后讨论如何辨识大语言模型的道德倾向，以及确保其符合伦理的复杂之处。本章还研究神经模型的毒性退化（neural toxic degeneration），介绍缓解和抑制这种倾向的策略。后续内容转向伦理中的隐私问题，特别是语言模型面对隐私攻击时的脆弱性，最后系统介绍针对大语言模型的隐私增强技术，着重说明保护用户数据与交互的先进方法。

\noindent\textbf{第 7 章\quad 检索增强生成}\enspace 本章阐述检索增强生成的基本要素，说明设计 RAG 系统时要考虑的问题。我们考察可纳入 RAG 工作流程的各种模块化改进，以扩展功能、减轻潜在漏洞的影响。随后分析评估 RAG 表现的主要指标，重点关注稠密检索（dense retrieval）过程的准确性及聊天机器人回答查询的效果。

\noindent\textbf{第 8 章\quad 生产环境中的大语言模型}\enspace 本章关注大语言模型在运行和工程上的问题，尤其关注已经逐步成为各类实际应用组成部分的提示方法。它系统介绍怎样在生产环境中有效部署大语言模型，讨论如何选择合适的模型、理解评估指标、开展基准评测，以及怎样从延迟、成本、质量、可适配性和维护等方面优化系统。章中给出所需的工具和技术，带领读者掌握大语言模型应用开发中的细节。

\noindent\textbf{第 9 章\quad 多模态大语言模型}\enspace 本章深入探讨迅速发展的多模态大语言模型，这一领域代表语言建模的一项重要进展。我们给出一个通用框架，从理论和实践两方面讨论各个组成部分，并与先进实现一一对应。本章考察如何把传统大语言模型的指令调优、上下文学习、思维链提示（chain-of-thought prompting）和对齐调优等技术用于多模态情境，说明这些改进如何提高模型跨模态适应和推理的能力。书中详述各类基准、数据集以及架构差异，重点介绍 Flamingo、Video-LLaMA 和 NExT-GPT 三种领先的模型，全面概述这些模型并说明它们在通用框架中的对应关系。

\noindent\textbf{第 10 章\quad 大语言模型的演进与新前沿}\enspace 全书最后一章回顾大语言模型的演进，重点梳理重要趋势和发展。它考察如何转向合成数据（synthetic data）以继续扩大模型规模，以及如何扩展上下文窗口（context window）以增强模型理解信息的能力。本章还讨论训练技术如何提高知识迁移（knowledge transfer）的效率、加深迁移程度，以及如何从传统 Transformer 架构转向状态空间模型（state space model）等具有更好可扩展性和效率的方法。本章也关注模型变小、效率提高、技术普及及领域专用模型兴起的趋势。这些趋势表明人工智能方案正变得更易定制、更易获取，也更贴合具体行业。此外，本章探讨大语言模型技术的前沿及其在基于智能体的应用和搜索引擎中的使用，这些应用正在逐步取代传统技术。

\frontsection{如何阅读本书}
本书提供灵活的阅读路径，以适应不同背景和兴趣的读者。例如，产业界数据科学家若着重关注如何用自定义数据集微调大语言模型、了解相关挑战及缓解措施，并把模型部署到生产环境，可以重点阅读第 1、3、4、7、8、9 章。

我们尽量安排章节，使较复杂的主题建立在较基础的概念之上。阅读时也请留意，材料有时并非严格按先后依赖关系编排。例如，第 2 章提到一些基础大语言模型，它们所依赖的训练技术要到第 6 章才得到充分解释。遇到这类情况，我们常给出书内参照，指向对有关信息讲得更细的章节。读者可以围绕感兴趣的主题，按需要前后跳读。开始下一章之前，也无需完全掌握此前各章。

每章的重要内容用灰底框突出显示，方便读者辨识并查找重要概念和信息。
\begin{keypoint}
这是一个重要概念。
\end{keypoint}

各章都设有带醒目标识的“实用提示”（Practical Tips）框，给出实际建议与策略。
\begin{practicaltip}
这些框介绍部署、调优、定制、工具和参数等方面的重要策略，给出将书中概念用于实际工作的具体做法。
\end{practicaltip}

书中的提示（prompt）统一采用标准清单样式。
\begin{lstlisting}[caption={GLUE 的 BoolQ 示例}]
passage: "Look What You Made Me Do" is a song recorded by
   American singer-songwriter Taylor Swift, released on August
   24, 2017 by Big Machine Records as the lead single from her
   sixth studio album Reputation (2017). Swift wrote the song
   with her producer Jack Antonoff. "Look What You Made Me Do"
   is an electroclash and pop song, with lyrics about various
   issues that built Swift's reputation. Right Said Fred band
   members Fred Fairbrass, Richard Fairbrass, and Rob Manzoli
   are also credited as songwriters, as it interpolates the
   melody of their song "I'm Too Sexy" (1991).

question: "did taylor swift write look what you made me do"

label: 1
\end{lstlisting}
\transnote{清单保留任务输入。段落说明 Taylor Swift 与制作人 Jack Antonoff 共同创作了歌曲，问题询问 Swift 是否创作了这首歌，标签 \texttt{1} 表示肯定。}

书中 Python 代码及清单采用清晰、统一的格式，便于读者理解和实践与主题有关的编程概念。
\begin{lstlisting}[language=Python,caption={初始化 BERT 词元切分器和模型的 Python 代码}]
from transformers import BertTokenizer, BertForMaskedLM
tokenizer = BertTokenizer.from_pretrained("bert-base-uncased")
model = BertForMaskedLM.from_pretrained("bert-base-uncased")
\end{lstlisting}

本书所有教程都托管在专用的 GitHub 仓库 \url{https://github.com/springer-llms-deep-dive}。仓库按章分设文件夹，包含 Jupyter 笔记本和相关代码。读者可在 Google Colab 中使用 GPU 运行设置，以取得更好的效率。我们建议订阅 Colab Pro。相对于增加的计算能力，我们认为订阅费用不高。我们有意让不同算力预算的从业者都能使用这些教程。不过，若不为 Colab Pro 或性能相当的 GPU 资源支付少量费用，部分笔记本可能遇到内存不足的问题。


% PDF 18: 致谢。
\frontchapter{致谢}
没有许多人的大力帮助，本书就无法成稿。首先感谢 Springer，特别是我们的编辑 Paul Drougas 和协调人 Jacob Shmulewitz。他们与我们密切合作，促成本书出版。感谢 Smarsh 给予我们处理现实世界中多模态、多语言难题的机会，也感谢公司培育的研究和创新氛围，它对本书的工作影响深远。

我们衷心感谢以下人士（按英文字母顺序排列）。Felipe Blanco、Shekar Gothoskar、Gaurav Harode、Sarang Kayande 博士、Ankit Mittal、Sasi Mudigonda、Raj Pai、Gokul Patel、Sachin Phadnis、Ross Turner 博士、Sameer Vajre 和 Vedant Vajre 对内容提出的反馈、建议以及他们作出的贡献，对完成本书至关重要。最后，我们诚挚感谢阅读、审阅本书并为序言和推荐语供稿的行业专家和研究人员。你们的见解和专业知识弥足珍贵。特别感谢 Ajit Jaokar、Shalini Govil Pai、Sherry Marcus 博士、Prithvi Prabhu、Amarda Shehu 博士和 Daniel Svonava。

% PDF 19: 声明。
\frontchapter{声明}
\noindent\textbf{利益冲突}\enspace 作者没有与本书内容相关的利益冲突需要申报。

\noindent\textbf{图表重制}\enspace 本书中据研究论文重制或改编的图均由作者绘制。有些图可能没有在图中明确写出原始出处，正文相应各节均恰当引用了原作。书中所有表格也遵循同样的做法。

\noindent\textbf{伦理审批}\enspace 本书不包含以人或动物为参与对象的一手研究，因此无须伦理审批。

% PDF 20--21: 符号表。
\fronttocchapter{符号表}
\begingroup
\small
\setlength{\LTpre}{0pt}
\renewcommand{\arraystretch}{1.25}
\begin{longtable}{@{}>{$}l<{$}>{\raggedright\arraybackslash}p{.70\linewidth}@{}}
\multicolumn{2}{@{}l}{\textbf{微积分}}\\
\approx & 近似相等\\
\lvert\symbf{A}\rvert & 矩阵 $\symbf{A}$ 的 $L_1$ 范数（norm）\\
\lVert\symbf{A}\rVert & 矩阵 $\symbf{A}$ 的 $L_2$ 范数\\
\frac{\mathrm{d}a}{\mathrm{d}b} & $a$ 对 $b$ 的导数\\
\frac{\partial a}{\partial b} & $a$ 对 $b$ 的偏导数\\
\nabla_{\symbf{x}}Y & $Y$ 对 $\symbf{x}$ 的梯度\\
\nabla_{\symbf{X}}Y & $Y$ 对 $\symbf{X}$ 的导数构成的矩阵\\[3pt]
\multicolumn{2}{@{}l}{\textbf{数据集}}\\
\mathcal{D} & 数据集，样例与相应目标的集合
$\{(\symbf{x}_1,y_1),(\symbf{x}_2,y_2),\ldots,(\symbf{x}_n,y_n)\}$\\
\mathcal{X} & 所有可能输入构成的空间\\
\mathcal{Y} & 所有可能输出构成的空间\\
y_i & 第 $i$ 个样例的目标标签\\
\hat{y}_i & 第 $i$ 个样例的预测标签\\
\mathcal{L} & 对数似然损失（log-likelihood loss）\\
\Omega & 学得的参数\\[3pt]
\multicolumn{2}{@{}l}{\textbf{函数}}\\
f\colon\mathbb{A}\to\mathbb{B} & 将集合 $\mathbb{A}$ 中的值映射到集合 $\mathbb{B}$ 的函数 $f$\\
f(\symbf{x};\theta) & 关于 $\symbf{x}$、以 $\theta$ 为参数的函数。为使记号简明，常简写为 $f(\symbf{x})$。\\
\log x & $x$ 的自然对数\\
\sigma(a) & 逻辑 S 形函数（logistic sigmoid），$\frac{1}{1+\exp(-a)}$\\
{[\![a\ne b]\!]} & 括号内条件为真时取 $1$、否则取 $0$ 的函数\\
{}_x f(x) & 使 $f(x)$ 最小的自变量的集合。
\( {}_x f(x)=\{x\mid f(x)=\min_{x'} f(x')\}\)\\
{}_x f(x) & 使 $f(x)$ 最大的自变量的集合。
\( {}_x f(x)=\{x\mid f(x)=\max_{x'} f(x')\}\)\\
\multicolumn{2}{@{}p{\linewidth}@{}}{\small \textbf{译注}\quad
原书第 xxiii 页最后两个函数记号缺少算子，右侧集合分别定义了 $\arg\min$ 与 $\arg\max$。}\\[3pt]
\multicolumn{2}{@{}l}{\textbf{线性代数}}\\
a & 标量（整数或实数）\\
\begin{bmatrix} a_1\\ \vdots\\ a_n\end{bmatrix} &
含有元素 $a_1$ 到 $a_n$ 的向量\\[3pt]
\begin{bmatrix} a_{1,1}&\cdots&a_{1,n}\\
\vdots&\ddots&\vdots\\ a_{m,1}&\cdots&a_{m,n}\end{bmatrix} &
有 $m$ 行、$n$ 列的矩阵\\[3pt]
A_{i,j} & 矩阵 $\symbf{A}$ 第 $i$ 行、第 $j$ 列的值\\
\symbf{a} & 向量（维数由上下文确定）\\
\symbf{A} & 矩阵（维数由上下文确定）\\
\symbf{A}^{\top} & 矩阵 $\symbf{A}$ 的转置\\
\symbf{A}^{-1} & 矩阵 $\symbf{A}$ 的逆\\
\symbf{I} & 单位矩阵（维数由上下文确定）\\
\symbf{A}\cdot\symbf{B} & 矩阵 $\symbf{A}$ 和 $\symbf{B}$ 的点积（dot product）\\
\symbf{A}\times\symbf{B} & 矩阵 $\symbf{A}$ 和 $\symbf{B}$ 的叉积（cross product）\\
\symbf{A}\circ\symbf{B} & 逐元素乘积（Hadamard product）\\
\symbf{A}\otimes\symbf{B} & 矩阵 $\symbf{A}$ 和 $\symbf{B}$ 的克罗内克积（Kronecker product）\\
\symbf{a};\symbf{b} & 向量 $\symbf{a}$ 与 $\symbf{b}$ 的拼接\\[3pt]
\multicolumn{2}{@{}l}{\textbf{概率}}\\
\mathbb{E} & 期望值\\
P(A) & 事件 $A$ 的概率\\
X\sim\mathcal{N}(\mu,\sigma^2) & 从均值为 $\mu$、方差为 $\sigma^2$ 的高斯（正态）分布中抽取的随机变量 $X$\\[3pt]
\multicolumn{2}{@{}l}{\textbf{集合}}\\
\mathbb{A} & 一个集合\\
\mathbb{R} & 实数集\\
\mathbb{C} & 复数集\\
\varnothing & 空集\\
\{a,b\} & 由元素 $a$ 和 $b$ 构成的集合\\
\{1,2,\ldots,n\} & 从 $1$ 到 $n$ 的所有整数组成的集合\\
\{a_1,a_2,\ldots,a_n\} & 含 $n$ 个元素的集合\\
a\in\mathbb{A} & 值 $a$ 是集合 $\mathbb{A}$ 的元素\\
{[a,b]} & 从 $a$ 到 $b$ 的实数构成的集合，包含 $a$ 和 $b$\\
{[a,b)} & 从 $a$ 到 $b$ 的实数构成的集合，包含 $a$，不包含 $b$\\
a_{1:m} & 元素集 $\{a_1,a_2,\ldots,a_m\}$（为便于书写）\\
\end{longtable}
\endgroup
除非另有说明，本书大多采用上面的记号。

% PDF 22--29: 目录随中文分页生成。
\clearpage
\pdfbookmark[0]{目录}{contents}
\tableofcontents
% PDF 30--31: 选列缩写。
\frontchapter{常用缩写}
\begingroup
\footnotesize
\renewcommand{\arraystretch}{1.00}
\begin{longtable}{@{}p{.14\linewidth}>{\raggedright\arraybackslash}p{.49\linewidth}>{\raggedright\arraybackslash}p{.29\linewidth}@{}}
\textbf{缩写} & \textbf{英文全称} & \textbf{中文释义}\\ \midrule
\endfirsthead
\textbf{缩写} & \textbf{英文全称} & \textbf{中文释义}\\ \midrule
\endhead
\abb{AI}{Artificial Intelligence}{人工智能}
\abb{BLEU}{Bilingual Evaluation Understudy}{双语评估替代指标}
\abb{CBOW}{Continuous Bag-of-Words}{连续词袋}
\abb{CBS}{Categorical Bias Score}{类别偏见分数}
\abb{CLIP}{Contrastive Language-Image Pre-training}{对比语言—图像预训练}
\abb{CNN}{Convolutional Neural Network}{卷积神经网络}
\abb{CoT}{Chain-of-Thought}{思维链}
\abb{CPT}{Continual Pre-Training}{持续预训练}
\abb{CUDA}{Compute Unified Device Architecture}{统一计算设备架构}
\abb{DMN}{Dynamic Memory Network}{动态记忆网络}
\abb{DPO}{Direct Preference Optimization}{直接偏好优化}
\abb{EOS}{End-of-Sentence (token)}{句末词元}
\abb{GeLU}{Gaussian Error Linear Unit}{高斯误差线性单元}
\abb{GLU}{Gated Linear Unit}{门控线性单元}
\abb{GNN}{Graph Neural Network}{图神经网络}
\abb{GPT}{Generative Pre-trained Transformer}{生成式预训练 Transformer}
\abb{GPU}{Graphical Processing Unit}{图形处理器}
\abb{HHH/3H}{Helpful, Honest, and Harmless}{有帮助、诚实且无害}
\abb{ICL}{In-Context Learning}{上下文学习}
\abb{IT}{Instruction Tuning}{指令调优}
\abb{ITG}{Image-Text Generation}{图像—文本生成}
\abb{ITM}{Image-Text Matching}{图像—文本匹配}
\abb{KD}{Knowledge Distillation}{知识蒸馏}
\abb{KL}{Kullback-Leibler}{库尔巴克—莱布勒}
\abb{LLM}{Large Language Model}{大语言模型}
\abb{LLMOps}{Large Language Model Operations}{大语言模型运维}
\abb{LM}{Language Model}{语言模型}
\abb{LoRA}{Low-Rank Adaptation}{低秩适配}
\abb{LPBS}{Log-Probability Bias Score}{对数概率偏见分数}
\abb{LSTM}{Long Short-Term Memory}{长短期记忆}
\abb{MDP}{Markov Decision Process}{马尔可夫决策过程}
\abb{ME}{Modality Encoder}{模态编码器}
\abb{MM-IT}{Multimodal Instruction Tuning}{多模态指令调优}
\abb{MLM}{Masked Language Modeling}{掩码语言建模}
\abb{MLOps}{Machine Learning Operations}{机器学习运维}
\abb{MM-COT}{Multimodal Chain-of-Thought}{多模态思维链}
\abb{MM-ICL}{Multimodal In-Context Learning}{多模态上下文学习}
\abb{MMLLM}{Multimodal Large Language Model}{多模态大语言模型}
\abb{MoE}{Mixture of Experts}{混合专家}
\abb{MRR}{Mean Reciprocal Rank}{平均倒数排名}
\abb{nDCG}{Normalized Discounted Cumulative Gain}{归一化折损累积增益}
\abb{NER}{Named-Entity Recognition}{命名实体识别}
\abb{NLG}{Natural Language Generation}{自然语言生成}
\abb{NLI}{Natural Language Inference}{自然语言推断}
\abb{NLP}{Natural Language Processing}{自然语言处理}
\abb{OCR}{Optical Character Recognition}{光学字符识别}
\abb{PEFT}{Parameter-Efficient Fine-Tuning}{参数高效微调}
\abb{PII}{Personally identifable information}{个人身份可识别信息}
\abb{PLM}{Pre-trained Language Model}{预训练语言模型}
\abb{PPO}{Proximal Policy Optimization}{近端策略优化}
\abb{PTFT}{Pre-Train and Fine-Tune}{预训练再微调}
\abb{PTQ}{Post-training Quantization}{训练后量化}
\abb{PTS}{Pre-training From Scratch}{从头预训练}
\abb{QA}{Question Answering}{问答}
\abb{QLoRA}{Quantized Low-Rank Adaptation}{量化低秩适配}
\abb{RAG}{Retrieval-Augmented Generation}{检索增强生成}
\abb{RL}{Reinforcement Learning}{强化学习}
\abb{RLHF}{Reinforcement Learning with Human Feedback}{基于人类反馈的强化学习}
\abb{RNN}{Recurrent Neural Network}{循环神经网络}
\abb{ROUGE}{Recall-Oriented Understudy for Gisting Evaluation}{面向召回的摘要评估替代指标}
\abb{RRF}{Reciprocal Rank Fusion}{倒数排名融合}
\abb{SFT}{Supervised Fine-Tuning}{监督微调}
\abb{SMoE}{Sparse Mixture of Experts Model}{稀疏混合专家模型}
\abb{SOTA}{State of the Art}{当前先进水平}
\abb{SQL}{Structured Query Language}{结构化查询语言}
\abb{SVD}{Singular Value Decomposition}{奇异值分解}
\abb{TCO}{Total Cost of Ownership}{总拥有成本}
\abb{TF-IDF}{Term Frequency/Inverse Document Frequency}{词频—逆文档频率}
\abb{TI}{Task Instructions}{任务指令}
\abb{ToT}{Tree-of-Thoughts}{思维树}
\abb{TPU}{Tensor Processing Unit}{张量处理单元}
\abb{TRPO}{Trust Region Policy Optimization}{信赖域策略优化}
\abb{VeRA}{Vector-Based Random Matrix Adaptation}{基于向量的随机矩阵适配}
\abb{VLM}{Visual Linguistic Matching}{视觉语言匹配}
\abb{VQA}{Visual Question Answering}{视觉问答}
\abb{VSM}{Video-Subtitle Matching}{视频—字幕匹配}
\end{longtable}
\endgroup

